% !TEX root = ../../../main.tex
% Block card: preprocessing.bandpass -- "Bandpass filter" (Part I, the blocks chapter).
\cardfield{What it does}
Keeps only the changes in the signal that happen at a middle range of speeds, removing both
the slow drift underneath and the fast jitter on top.

\cardfield{When to use it}
A fungal recording mixes movements at very different speeds: a level that wanders over hours,
events that rise and fall over minutes, and fast noise from second to second. Think of it as a
song with a deep hum, a melody and hiss: a band-pass filter is a tone control that keeps the
melody and turns down both the hum and the hiss. Use it when the events you care about live at
one time-scale and you want a detector to see only that. When Discovery applies a template
over a band scope, it puts this block at the head of the chain once per band, taking the edges
from the band list on the Settings page (Analysis defaults).

\cardfield{What goes in, what comes out}
In goes a span of one channel. Out comes a signal of exactly the same length and sampling rate,
now centred on zero, holding only the middle speeds.

\cardfield{Settings that matter}
\emph{low\_hz} and \emph{high\_hz} are the edges of the band, in cycles per second (Hz). A
movement that takes 100 seconds to repeat is 0.01\,Hz; one that takes 10 seconds is 0.1\,Hz,
so the defaults keep what changes over roughly 10 to 100 seconds. \emph{order} sets how sharp
the edges are: higher cuts more steeply but rings more.

\cardfield{Watch out for}
\emph{high\_hz} must be below half the sampling rate (0.5\,Hz on a 1\,Hz recording), or the
block stops with an error. A single missing sample anywhere in the span turns the whole output
into missing values. A sudden drop comes out with ripples on both sides of it, including
\emph{before} it starts, and its depth changes: the filter reshapes events, so measure their
size on the unfiltered trace.
